\documentclass[twocolumn]{aastex631}

\usepackage{amsmath}
\usepackage{multirow}
\usepackage{natbib}
\usepackage{graphicx} 
\usepackage{aas_macros}

\begin{document}

The search for stable and observationally viable extensions to General Relativity, particularly within Horndeski theories, is crucial for understanding dark energy and modified gravity, yet these theories often suffer from theoretical instabilities. This paper investigates whether a novel "Disformal-Generalized Galilean Co-invariance," characterized by the simultaneous invariance of the action under field-dependent disformal metric transformations and a generalized Galilean symmetry for the scalar field, can protect a subclass of Horndeski theories from both classical and one-loop quantum instabilities. Through rigorous analytical derivation, we identified the unique Horndeski subclass satisfying this co-invariance, revealing it to be precisely covariant Galileon theories with constant coefficients for the quartic and quintic terms, while allowing for arbitrary k-essence and a $\phi$-dependent cubic Galileon term. Our classical stability analysis demonstrated that the tensor sector of this co-invariant subclass requires a negative kinetic term coefficient ($c_4 < 0$) for ghost-freedom. Critically, we showed that this underlying symmetry structure inherently ensures the non-renormalization of ghost modes at the one-loop level, affirming its theoretical robustness against quantum corrections. However, confronting the model with stringent observational constraints, specifically the speed of gravitational waves ($c_T=1$) from GW170817 in the late universe, leads to a direct contradiction. In a de Sitter background, this observation forces the very parameter governing the tensor kinetic term ($c_4$) to be zero, which implies a pathological theory with a vanishing kinetic term, violating the classical no-ghost condition. This work thus presents a no-go theorem for Disformal-Generalized Galilean Co-invariant Horndeski theories, demonstrating a fundamental tension between the symmetries required for theoretical stability and the demands of current cosmological observations.

\end{document}
                